%! Simplifying Radicals; Adding \& Subtracting — Unit 5, Lesson 5.2 — Exit Ticket (Answer Key)
\documentclass[10pt]{article}
\usepackage{algebra2-article}
\usepackage{algebra2-key}   % pulls in -boxes; do NOT also load -boxes
\newcommand{\TallMath}[1]{$\displaystyle #1\rule[-1.4em]{0pt}{3.2em}$}

\begin{document}

\pageheader{Unit 5, Lesson 5.2}{Exit Ticket --- Answer Key}
\namedateperiod

\vspace{0.10in}

No notes, no calculator. Assume all variables represent positive real numbers.

\begin{enumerate}[leftmargin=1.8em, itemsep=10pt]
  \item \textbf{Simplify completely.}

        \vspace{4pt}
        (a) $\sqrt{75} = $ \ans{$5\sqrt{3}$} \hspace{0.5in}
        (b) $\sqrt[3]{54} = $ \ans{$3\sqrt[3]{2}$}

        \vspace{4pt}
        (c) $\sqrt{18x^{5}} = $ \ans{$3x^{2}\sqrt{2x}$} \hspace{0.4in}
        (d) $2\sqrt{12}+\sqrt{27} = $ \ans{$7\sqrt{3}$}

        \vspace{2pt}
        {\small\emph{Work:} (a) $75=25\cdot3$. \; (b) $54=27\cdot2$ (group in \emph{threes}). \;
        (c) $18=9\cdot2$ and $x^{5}=x^{4}\cdot x$. \;
        (d) $2\sqrt{12}=2\cdot2\sqrt3=4\sqrt3$ and $\sqrt{27}=3\sqrt3$, so $4\sqrt3+3\sqrt3=7\sqrt3$
        --- neither term was a like radical until it was simplified.}

  \item \textbf{Explain.} A classmate writes $\sqrt{9}+\sqrt{16}=\sqrt{25}=5$. Show with numbers
        that this is wrong, then state the rule that \emph{does} hold and say what operation it
        requires.
        \ansline{$\sqrt9+\sqrt{16}=3+4=7$, not $5$. So $\sqrt{a}+\sqrt{b}\ne\sqrt{a+b}$ --- a
        radical does not split over addition.\\[3pt]
        What does hold is the \textbf{product} property: $\sqrt{ab}=\sqrt a\cdot\sqrt b$. It
        requires \emph{multiplication} under the radical, not addition --- $\sqrt{9\cdot16}
        =3\cdot4=12$ checks out.}

  \item \textbf{SOL-style multiple choice.} Which expression is equivalent to
        $\sqrt{50}+\sqrt{32}$?
        \begin{enumerate}[label=\Alph*., leftmargin=1.9em, itemsep=1pt]
          \item $\sqrt{82}$
          \item $9\sqrt{2}$ \quad \textcolor{keyred}{\textbf{$\leftarrow$ correct}}
          \item $9\sqrt{4}$
          \item $20\sqrt{2}$
        \end{enumerate}
        \vspace{2pt}
        \begin{itemize}[leftmargin=1.4em, nosep]
          \item \textbf{B} is correct: $\sqrt{50}=5\sqrt2$ and $\sqrt{32}=4\sqrt2$, so $5\sqrt2+4\sqrt2=9\sqrt2$.
          \item \textbf{A} adds the radicands --- the error item 2 is about.
          \item \textbf{C} adds the coefficients correctly but then adds the radicands too.
          \item \textbf{D} multiplies the coefficients ($5\cdot4$) instead of adding them.
        \end{itemize}
\end{enumerate}

\begin{teachernote}
Graded for completion --- mistakes happen, blanks don't. Sort the stack into four piles to decide
how Lesson 5.3 opens:
\begin{itemize}[leftmargin=1.3em, nosep]
  \item \textbf{Got it} --- 1(a)--(d) right and item 3 = B.
  \item \textbf{Stopped early} --- $\sqrt{75}$ left as $5\sqrt3$ but $\sqrt{18x^5}$ answered
        $3x^2\sqrt{2x}$ only partly, or 1(b) answered $\sqrt[3]{54}=2\sqrt[3]{27}$-style
        pair-grouping. These students used a factor that was not the largest, or grouped in twos
        under a cube root. One prompt fixes it: ``Is there still a perfect cube hiding in there?''
  \item \textbf{Adders} --- item 3 = A or C, or item 2 unanswered. This is the pile that matters:
        the belief that a radical distributes over addition survives everything except being
        checked numerically, which is exactly what item 2 asks for.
  \item \textbf{Did not simplify first} --- 1(d) answered ``cannot combine.'' They applied the
        like-radical rule honestly to $\sqrt{12}$ and $\sqrt{27}$ and concluded, reasonably, that
        the radicands differ. The lesson's central habit --- \emph{simplify, then decide} --- has
        not landed yet, and 5.3's distributing work will magnify it.
\end{itemize}
Item 1(c) is the only algebraic radicand; if the errors cluster there, re-run the exponent-division
picture ($x^5=x^4\cdot x$, or $x^{5/2}=x^2\cdot x^{1/2}$) rather than re-teaching the numeric case.
\end{teachernote}

\end{document}
